\chapter{Semisimple Modules and Rings}

In this chapter, when we speak of a module (without further specification), we mean a left module. Let A be a ring and M an A-module. For any \(a \in A\) , we denote the homothety \(x \mapsto ax\) of M by \(a_{M}\) . The mapping \(a \mapsto a_{M}\) is a homomorphism from the ring A to a subring \(A_{M}\) of \(\operatorname{End}_{\mathbf{Z}}(\mathrm{M})\) , which we call the ring of homotheties of M.

Unless stated otherwise, the algebras we consider are associative and unital; by a subalgebra of an algebra E, we mean a subalgebra containing the unit element of E; the algebra homomorphisms are assumed to be unital.

Let K be a commutative ring and L a commutative K-algebra. For any K-module E, we denote the L-module \(L \otimes_{K} E\) derived from it by extension of scalars by \(\mathrm{E}_{(\mathrm{L})}\) (II, §5, No.~1, p.~277). If A is a K-algebra and M a left A-module, then \(\mathrm{A}_{(\mathrm{L})}\) is endowed with a natural L-algebra structure (III, §1, No.~5, p.~433) and \(\mathrm{M}_{(\mathrm{L})}\) is endowed with the left \(\mathrm{A}_{(\mathrm{L})}\) -module structure with law of action given by the formula \((\lambda \otimes a)(\mu \otimes m) = \lambda \mu \otimes am\) .

For any commutative ring K and any group G, we denote the algebra \(\mathrm{K}^{(\mathrm{G})}\) of G over the ring K by K{[}G{]} (III, §2, No.~6, p.~446).

